\title{xLSTM: Extended Long Short-Term Memory}

\begin{abstract}
The introduction of the constant error carousel and gating mechanisms in the 1990s established the foundational principles of Long Short-Term Memory (LSTM), which have since become pivotal in the development of deep learning, most notably powering the earliest Large Language Models (LLMs). Despite LSTM’s enduring relevance, the rise of Transformer architectures—driven by inherently parallel self-attention mechanisms—has largely eclipsed LSTM models at greater scale. This paper revisits the LSTM framework with a central inquiry: To what extent can large-scale LSTMs, enhanced with innovations from contemporary LLM practices and refined to address historical shortcomings, rival recent models in language modeling? Our contributions are twofold: We propose exponential gating along with normalization and stabilization strategies, and we alter the canonical memory configuration of the LSTM. The resulting architectures include: (i) sLSTM, characterized by scalar-valued memory and updates, as well as novel approaches to memory mixing, and (ii) mLSTM, which features fully parallel operations through matrix-valued memory with a covariance-based update scheme. We integrate these upgraded LSTM variants into residual network structures, creating xLSTM modules that are then recursively composed to yield complete xLSTM architectures. The synergistic effect of exponential gating and reimagined memory organization significantly enhances the performance and scalability of xLSTMs, allowing them to achieve results competitive with current Transformers and State Space Models.\\
Code available at: \url{https://github.com/NX-AI/xlstm}
\end{abstract}

\section{Introduction}
\label{sec:intro}

The concepts behind Long Short-Term Memory (LSTM)~\citep{Hochreiter:91,Hochreiter:97nips,Hochreiter:97}, notably the constant error carousel and the gating mechanisms, were proposed as solutions to the vanishing gradient issue encountered in recurrent neural networks~\citep{Hochreiter:91,Hochreiter:00book}:

\begin{align}
\label{eq:lstm_idea}
  c_t \ = \ f_t \ c_{t-1} \ + \ 
          i_t \  z_t \ , \quad  
          h_t \ = \ o_t \ \psi({c_t}) \ . 
\end{align}

The constant error carousel functions by incrementally updating the cell state $c_{t-1}$ (depicted in green) using the cell inputs $z_t$, with the process regulated through sigmoid gates (shown in blue). The update is specifically modulated by the input gate $i_t$ and forget gate $f_t$, whereas the memory cell’s output—corresponding to the hidden state $h_t$—is controlled by the output gate $o_t$. Before contributing to the hidden state, the cell state undergoes normalization or squashing via the function $\psi$, after which output gating determines the final hidden state.

LSTMs have achieved notable success across a wide range of fields \citep{Hochreiter:01,Hochreiter:07,Schmidhuber:15}, maintaining dominance in text generation tasks prior to the introduction of Transformers in 2017~\citep{Vaswani:17}. Their capabilities have been extensively validated in tasks involving sequences, including text production~\citep{Graves:13,Karpathy:15blog}, handwriting synthesis~\citep{Graves:13}, machine translation using sequence-to-sequence models~\citep{Sutskever:14nips}, program analysis~\citep{Zaremba:14arxiv}, image captioning~\citep{Karpathy:15,Hossain:19}, code generation~\citep{Karpathy:15blog}, rainfall-runoff prediction~\citep{Kratzert:18,Kratzert:19}, and hydrological forecasting for flood warnings~\citep{Nearing:24}. In reinforcement learning scenarios, LSTMs stand as the leading sequence modeling architecture, as exemplified by their use in the AlphaStar system for StarCraft II~\citep{Vinyals:17short}, the OpenAI Five project for Dota 2~\citep{OpenAI:19}, and nuclear fusion’s magnetic control models~\citep{Degrave:22short}. Their strength lies in abstracting complex concepts—efficiently capturing semantic representations within their memory cells~\citep{Karpathy:15blog}—which is observable through constructs such as number and syntax neurons~\citep{Lakretz:19}, linguistic neurons~\citep{Bau:19}, and sentiment neurons~\citep{Radford:17arxiv}. Despite emerging alternatives, LSTMs remain integral to mission-critical systems~\citep{Degrave:22short,Nearing:24} and have demonstrated enduring relevance.

Although LSTMs have achieved significant breakthroughs, they present three primary shortcomings: (i) Irreversibility in storage updates. This shortfall is illustrated using the {\it Nearest Neighbor Search} scenario (refer also to Appendix~\ref{sec:appNearestNeighborSearch}): Given a reference vector, the model must traverse the entire sequence to locate the most similar vector, so that its corresponding value can be reported at the sequence’s end. As shown in the left portion of Figure~\ref{fig:lstmProblems}, the mean squared error for this challenge highlights how LSTMs have difficulty overwriting a previously stored item when a better candidate is subsequently found. In contrast, our proposed xLSTM addresses this deficit by employing exponential gating mechanisms. (ii) Constrained memory size—information must be compressed into single-dimensional cell states. We use the {\it Rare Token Prediction} task to demonstrate this: The right portion of Figure~\ref{fig:lstmProblems} presents perplexity scores for token prediction on Wikitext-103~\citep{Merity:17}, separated by token frequency. The LSTM exhibits diminished performance, particularly with rare tokens, as a consequence of its limited memory representation. Our xLSTM model resolves this issue through its use of matrix-based memory. (iii) Inability to parallelize processing, stemming from entanglement of hidden states across consecutive time steps (hidden-hidden connectivity), thereby necessitating strictly sequential computation.

The shortcomings inherent to LSTM models have led to the development and widespread adoption of Transformers~\citep{Vaswani:17} for language modeling tasks. This raises a pertinent question: if we address these LSTM limitations and expand LSTMs to match the scale of present-day Large Language Models, what levels of performance might be attainable in language modeling?

\section{Extended Long Short-Term Memory}
\label{sec:xLSTM}

To address the constraints inherent in LSTM, the Extended Long Short-Term Memory (xLSTM) framework presents two principal changes to the foundational concept of Equation~\eqref{eq:lstm_idea}. Specifically, xLSTM integrates exponential gating mechanisms and introduces innovative memory architectures, thereby giving rise to two new variants within the LSTM taxonomy: (i) the sLSTM (refer to Section~\ref{sec:sLSTM}), which operates using scalar memory, scalar updates, and incorporates memory mixing, and (ii) the mLSTM (see Section~\ref{sec:mLSTM}), characterized by a matrix-based memory structure with covariance (outer product) update dynamics that facilitate complete parallelization. The utilization of exponential gating in both sLSTM and mLSTM constitutes a significant enhancement over the traditional LSTM model. In favor of parallel computation, the mLSTM forgoes memory mixing—removing the hidden-to-hidden recurrent pathways. Both architectures allow for generalization to scenarios involving multiple memory cells; in this setup, the sLSTM incorporates mixing between memory cells. Additionally, sLSTM can be designed with multiple heads, omitting inter-head memory mixing but maintaining memory interactions among cells within each head, a setup that synergizes with exponential gating to produce a novel form of memory mixing. In the case of mLSTM, configurations employing multiple heads and multiple cells serve equivalent roles.

By embedding these novel LSTM variants within residual block structures, we construct xLSTM blocks (refer to Section~\ref{sec:xLSTMblock}). Building architectures by stacking these xLSTM blocks with residual connections yields xLSTM architectures (refer to Section~\ref{sec:xLSTMarchitecure}). The overall composition and elements of the xLSTM architecture are illustrated in Figure~\ref{fig:xlstm_sketch}.

\subsection{Review of the Long Short-Term Memory}
\label{sec:orgLSTM}

The foundational concept behind LSTM~\citep{Hochreiter:91,Hochreiter:97nips,Hochreiter:97} was the proposal of a scalar memory cell serving as the principal component for information processing and retention, thereby addressing the problem of vanishing gradients~\citep{Hochreiter:91,Hochreiter:00book} through what is termed the constant error carousel, or the iterative cell state update. This memory cell structure integrates three distinct gating mechanisms: the input gate, the output gate, and the forget gate—the latter introduced by \citet{Gers:00}. The LSTM memory cell update equations at time step $t$ are given by:

\begin{align}
c_t \ &= \  f_t \ c_{t-1} \ + \ i_t \ z_t &  & &\text{cell state} \\
h_t  \ &= \ o_t \ \tilde{h}_t \ , & \tilde{h}_t \ &= \ \psi \left( c_t\right)
  &\text{hidden state} \\
z_t \ &= \ \varphi \left( \tilde{z}_t \right) \ , 
  &\tilde{z}_t \ &=  \ \mathbf{w}^\top_{z} \ \mathbf{x}_t \ + \
  r_{z}  h_{t-1} \ + \  b_{z} \ \
  &\text{cell input} \\
i_t \ &= \ \sigma \left( \tilde{i}_t  \right) \ , 
  &\tilde{i}_t \ &= \ \mathbf{w}^\top_{i} \ \mathbf{x}_t \ + \
  r_{i}  \ h_{t-1} \ + \  b_{i} \ \
  &\text{input gate} \\
f_t \ &= \ \sigma \left(  \tilde{f}_t \right) \ , 
  &\tilde{f}_t \ &= \ \mathbf{w}^\top_{f} \ \mathbf{x}_t  \ + \
  r_{f}  \ h_{t-1} \ + \  b_{f} \ \
  &\text{forget gate} \\
o_t \ &= \ \sigma \left( \tilde{o}_t \right) \ , 
  &\tilde{o}_t  \ &= \ \mathbf{w}^\top_{o} \ \mathbf{x}_t \ + \
  r_{o}  \ h_{t-1} \ + \  b_{o} \ \
  &\text{output gate} 
\end{align}

The input weight vectors $\mathbf{w}_{z}$, $\mathbf{w}_{i}$,
$\mathbf{w}_{f}$, and $\mathbf{w}_{o}$ represent the weights from the external input $\mathbf{x}_t$ to the cell input as well as to the input, forget, and output gates, respectively. Meanwhile, the recurrent weights $r_{z}$, $r_{i}$, $r_{f}$, and $r_{o}$ denote the connections from the previous hidden state $h_{t-1}$ to the cell input and each respective gate (input, forget, and output). The terms $b_{z}$, $b_{i}$, $b_{f}$, and $b_{o}$ serve as biases for these elements. The activation functions $\varphi$ and $\psi$ typically take the form of $\tanh$, where $\varphi$ applies to the cell input and $\psi$ to the hidden state. Notably, $\psi$ is utilized to restrict the range of the cell state, which otherwise could diverge without such a transformation. Sigmoid functions are employed for all gates, formally given by $\sigma \left( x \right)= 1/(1+\exp(-x))$. In subsequent model versions, collections of scalar memory cells $\mathbf{c}_t \in \mathbb{R}^{d}$ are stacked into a vector, which enables recurrent weight matrices $\mathbf{R} \in \mathbb{R}^{d \times d}$ to interconnect outputs of different memory cells within the architecture~\citep{Greff:15}. Further information is presented in Appendix~\ref{sec:apporgLSTM}. Investigation via ablation has indicated that every part of the memory cell structure plays an essential role in its overall function~\citep{Greff:15}.

\subsection{sLSTM}
\label{sec:sLSTM}

To enable LSTMs to adjust their memory storage dynamically, we incorporate exponential gates (shown in red) alongside mechanisms for normalization and stabilization. Specifically, the input and forget gates are modified to employ exponential activation functions. For the purpose of normalization, a normalizer state is implemented, which accumulates the sum of the input gate multiplied by all subsequent forget gates.

The scalar sLSTM forward pass is:

\begin{align}
\label{eq:slstmforward}
c_t \ &= \  f_t \ c_{t-1} \ + \ i_t \ z_t & & 
 &\text{cell state} \\
n_t \ &= \  f_t \ n_{t-1} \ + \ i_t  & & 
 &\text{normalizer state} \\
h_t \ &= \ o_t \ \tilde{h}_t \ ,
  &\tilde{h}_t \ &= \ c_t / n_t
&\text{hidden state} \\
z_t \ &= \ \varphi \left( \tilde{z}_t \right) \ , 
  &\tilde{z}_t \ &=  \ \mathbf{w}^\top_{z} \ \mathbf{x}_t \ + \
  r_{z}  h_{t-1} \ + \  b_{z}
  &\text{cell input} \\
i_t \ &= \ \!\exp \left( \tilde{i}_t  \right) \ , 
  &\tilde{i}_t \ &= \ \mathbf{w}^\top_{i} \ \mathbf{x}_t \ + \
  r_{i}  \ h_{t-1} \ + \  b_{i}
  &\text{input gate} \\
f_t \ &= \ \sigma \left(  \tilde{f}_t \right) \ \text{OR} \
\!\exp \left(  \tilde{f}_t \right) \ ,
  &\tilde{f}_t \ &= \ \mathbf{w}^\top_{f} \ \mathbf{x}_t  \ + \
  r_{f}  \ h_{t-1} \ + \  b_{f}
  &\text{forget gate} \\
o_t \ &= \ \sigma \left( \tilde{o}_t \right) \ , 
  &\tilde{o}_t  \ &= \ \mathbf{w}^\top_{o} \ \mathbf{x}_t \ + \
  r_{o}  \ h_{t-1} \ + \  b_{o}
  &\text{output gate} 
\end{align}

We adapt the foundational LSTM gating mechanisms—specifically, gating based on input and/or hidden state along with a bias term—for integration into the proposed architectures. Since exponential activations may generate excessively large outputs and risk overflow, we mitigate this issue by employing an auxiliary state, denoted as \mbox{$m_t$~\citep{Milakov:18arxiv}}, to stabilize the gating process:

\begin{align}
\label{eq:slstmstabil}
m_t \ &= \ \max \big( \log(f_t) + m_{t-1}, \log(i_t) \big) &\text{stabilizer state} \\
i'_t \ &= \ \!\exp \left( \log(i_t) - m_t \right) \ = \ \!\exp \left( \tilde{i}_t - m_t \right) \ \, 
  &\text{stabil. input gate} \\
f'_t \ &= \ \!\exp \left( \log(f_t) + m_{t-1} - m_t \right) \ \, 
  &\text{stabil. forget gate}
\end{align}

In Appendix~\ref{sec:appsLSTM}, we demonstrate that substituting $f_t$ with $f'_t$ and $i_t$ with $i'_t$ in the forward computation has no effect on either the network's final output or the gradients of the loss function concerning the parameters.

\paragraph{New Memory Mixing.} Similar to the original LSTM (refer to Appendix~\ref{sec:appsLSTM}), sLSTM is capable of supporting several memory cells. The use of multiple memory cells facilitates the process of memory mixing, achieved through recurrent connections $\mathbf{R}_{\mathbf{z}}$, $\mathbf{R}_{\mathbf{i}}$, $\mathbf{R}_{\mathbf{f}}$, and $\mathbf{R}_{\mathbf{o}}$ linking the hidden state vector $\mathbf{h}$ to both the memory cell input $\mathbf{z}$ as well as the gates $\mathbf{i}$, $\mathbf{f}$, and $\mathbf{o}$. A distinctive feature introduced in this context is exponential gating, which uniquely influences the memory mixing dynamics. In this enhanced sLSTM, several heads can be utilized, allowing memory mixing to occur independently within each head, but not between different heads. The combined use of heads and exponential gating within sLSTM thus provides a novel approach to memory mixing.

\subsection{mLSTM}
\label{sec:mLSTM}

To improve the storage capability of LSTMs, we generalize the memory cell from a scalar value $c \in \mathbb{R}$ to a matrix form $\mathbf{C} \in \mathbb{R}^{d \times d}$. With this adjustment, information retrieval operates through matrix multiplication. At each time step $t$, a key-value pair consisting of the key $\mathbf{k}_t \in \mathbb{R}^d$ and the value $\mathbf{v}_t \in \mathbb{R}^d$—terminology aligned with that used in the Transformer model—are to be stored. Subsequently, at time $t+\tau$, one aims to recover the value $\mathbf{v}_t$ using a query vector $\mathbf{q}_{t+\tau} \in \mathbb{R}^d$. This approach corresponds to the framework described by Bidirectional Associative Memories (BAMs) \citep{Kohonen:72,Anderson:72,Nakano:72,Anderson:77}.

The covariance update rule~\citep{Sejnowski:77,Dayan:91} for storing a key-value pair is

\begin{align}
\mathbf{C}_t \ &= \ \mathbf{C}_{t-1} \ + \ \mathbf{v}_{t} \ \mathbf{k}_t^\top \ .
\end{align}

We apply layer normalization prior to mapping the inputs onto key and value representations, ensuring these components maintain zero mean. Employing the covariance update rule achieves optimal separation among retrieved binary vectors, which in turn maximizes the signal-to-noise ratio~\citep{Dayan:91}. Enhanced separability becomes feasible if retrieval is limited to pairwise comparisons, although this involves accepting the quadratic computational demands typical of attention mechanisms~\citep{Krotov:16,Krotov:17,Ramsauer:21}. Notably, the covariance update rule corresponds to the mechanism used in Fast Weight Programmers \citep{Schmidhuber:92ncfastweights,Schlag:21}, which have subsequently incorporated a fixed decay factor applied to $\mathbf{C}_{t-1}$ and a constant learning rate that scales $\mathbf{v}_{t} \mathbf{k}_t ^\top$~\citep{Ba:16}. Consistent with these insights, we incorporate the covariance update strategy into the LSTM architecture, where the forget gate functions as a decay parameter, the input gate acts analogously to a learning rate, and the output gate adjusts the amplitude of the vector that is retrieved.

In the context of this matrix memory, the normalizer state is constructed by summing key vectors, each scaled by the input gate as well as all subsequent forget gates. This approach allows the normalizer state to capture the magnitude of the gate activations over time. Since the inner product between the query and the normalizer state may approach zero, we adopt the practice from earlier work~\citep{Sun:23arxiv}, taking the absolute value of this inner product and imposing a minimum threshold (commonly set at 1.0) to maintain numerical stability. The forward computation for mLSTM is as follows:

\begin{align}
\label{eq:mlstm_recurrent_begin}
\mathbf{C}_t \ &= \  f_t \ \mathbf{C}_{t-1} \ + \ 
  i_t \ \mathbf{v}_t \ \mathbf{k}_t^\top & &  &\text{cell state} \\
\mathbf{n}_t \ &= \  f_t \ \mathbf{n}_{t-1} \ + \ 
  i_t \ \mathbf{k}_t & &  &\text{normalizer state} \\
\mathbf{h}_t  \ &= \ \mathbf{o}_t \ \odot \ \tilde{\mathbf{h}}_t
\ , \qquad \qquad 
\tilde{\mathbf{h}}_t \ = \   \mathbf{C}_t \mathbf{q}_t \ / \ 
  \max \left\{ |\mathbf{n}_t^\top \mathbf{q}_t|, 1 \right\} 
   &&&\text{hidden state} \\
\mathbf{q}_t \ &= \ \mathbf{W}_q \ \mathbf{x}_t \ + \ \mathbf{b}_q  & &  &\text{query input} \\
\mathbf{k}_t \ &= \ \frac{1}{\sqrt{d}} \mathbf{W}_k \ \mathbf{x}_t \ + \ \mathbf{b}_k  & &  &\text{key input} \\
\mathbf{v}_t \ &= \ \mathbf{W}_v \ \mathbf{x}_t \ + \ \mathbf{b}_v  & &  &\text{value input} \\
i_t \ &= \ \!\exp \!  \! \left( \tilde{i}_t  \right) \ , 
 \qquad \qquad \ \ \ \ \,
  \tilde{i}_t \ = \ \mathbf{w}^\top_{i} \ \mathbf{x}_t \ + \  b_{i} \ \ \ 
  &&&\text{input gate} \\
f_t \ &= \ \sigma \!  \left(  \tilde{f}_t \right) \ \text{OR} \ \!\exp \!  \! \left(  \tilde{f}_t \right) , \ \ \: \!
\tilde{f}_t \ = \ \mathbf{w}^\top_{f} \ \mathbf{x}_t  \ + \
  b_{f}
  &&& \text{forget gate} \\
\label{eq:mlstm_recurrent_end}
\mathbf{o}_t \ &= \ \sigma \left( \tilde{\mathbf{o}}_t \right) \ , \qquad \qquad \quad \ \ \ \,
  \tilde{\mathbf{o}}_t  \ = \ \mathbf{W}_{\mathbf{o}} \ \mathbf{x}_t \ + \
  \mathbf{b}_{\mathbf{o}}
  &&&\text{output gate} 
\end{align}

The mLSTM architecture is capable of incorporating several memory cells, analogous to the structure of standard LSTM networks. Within mLSTM, the use of multiple heads is functionally identical to using multiple memory cells, as memory intermixing does not occur. To ensure stability in the exponential gating mechanism of mLSTM, we apply the same stabilization procedures utilized for sLSTM, as outlined in Equation~\ref{eq:slstmstabil}. 
Due to the absence of memory blending in mLSTM, its recurrence relations can be re-expressed in a parallelized format. Additional information on this formulation can be found in Appendix~\ref{sec:appmLSTM}.

\subsection{xLSTM Architecture}
\label{sec:xLSTMarch}

\paragraph{xLSTM Blocks.}
\label{sec:xLSTMblock}
An xLSTM block is designed to non-linearly encode prior information within a high-dimensional representation, thereby improving the model’s ability to distinguish among different historical contexts. Distinguishing between histories is essential for making accurate predictions of subsequent sequence elements, such as the next token. This approach is motivated by Cover’s Theorem~\citep{Cover:65}, which asserts that patterns which are non-linearly embedded into higher-dimensional spaces have a greater chance of being separated via linear methods than they would in their original dimensionality. Two types of residual block configurations are explored: (i) A residual block employing post up-projection (akin to the structure in Transformers), where a non-linear function first summarizes historical information in the original feature space, followed by a linear projection into a high-dimensional space, a non-linear activation, and a subsequent linear projection back to the initial space; refer to the left side of Figure~\ref{fig:Backbones}
and the third column in Figure~\ref{fig:xlstm_sketch} for illustration.
A comprehensive schematic can be found in Figure~\ref{fig:appsLSTM_detailed}
in the appendix. (ii) A residual block utilizing pre up-projection (analogous to State Space Models), in which a linear mapping to a higher-dimensional space is first performed,

It captures information from previous timesteps in a non-linear fashion within a high-dimensional space, and subsequently applies a linear transformation to return to the original dimensionality. When an xLSTM block incorporates an sLSTM, we employ the post up-projection block. Conversely, for xLSTM blocks with an mLSTM, we implement the pre up-projection block, taking into account the increased memory capacity present in the high-dimensional space. For further illustration, please consult the left panel of Figure~\ref{fig:Backbones}, the third column of Figure~\ref{fig:xlstm_sketch}, or Figure~\ref{fig:appsLSTM_detailed} located in the appendix.

\paragraph{xLSTM Architecture. }
\label{sec:xLSTMarchitecure}
The xLSTM model is formed by arranging its components in a residual stack~\citep{Srivastava:15,He:16}. Our implementation adopts the widely employed pre-LayerNorm~\citep{Ba:16b} residual structure typical of state-of-the-art Large Language Models. This design is illustrated in the final column of Figure~\ref{fig:xlstm_sketch}.

\subsection{Memory and Speed Considerations}

In contrast to Transformers, xLSTM architectures exhibit linear computational requirements and maintain fixed memory usage regardless of sequence length. Owing to the compressive nature of xLSTM memory, these networks are particularly appropriate for industrial settings and edge device deployment.

mLSTM’s memory design eliminates the need for extra parameters, but incurs substantial computational costs because both its internal memory and update mechanism are represented as $d \times d$ matrices. This approach prioritizes memory capacity while increasing computational demands. Despite this, the heavy computations can be efficiently parallelized on GPUs, which means their impact on actual runtime remains minimal.

Although mLSTM can be parallelized similarly to FlashAttention~\citep{Dao:22,Dao:23} or GLA~\citep{Yang:23arxiv}, sLSTM lacks this parallelization capability because of its memory mixing (hidden-hidden connections). Nonetheless, we have created an efficient CUDA implementation that incorporates GPU memory optimizations down to the register level, resulting in performance that is generally less than twice as slow as mLSTM.

\section{Related Work}
\label{sec:related}

\paragraph{Linear Attention.} Various approaches have been introduced to address the quadratic scaling of Transformer attention with respect to input sequence length, seeking to reduce the computational cost to linear. For instance, the Synthesizer generates learned synthetic attention weights without requiring direct interactions between tokens~\citep{Tay:20arxiv}. Linformer achieves a linear complexity by modeling self-attention through a low-rank projection, thereby approximating the attention operation linearly~\citep{Wang:20arxiv}. The Linear Transformer adapts the attention computation itself into a linear form~\citep{Katharopoulos:20}. Performer approximates the softmax function in attention using the positive orthogonal random features technique, yielding a linear complexity~\citep{Choromanski:21}. Additionally, some recent architectures, such as Structured Global Convolution (SGConv)~\citep{Li:22} and the Hyena Hierarchy~\citep{Poli:23}, substitute conventional attention with fast, long-range convolutional operations.

\paragraph{State Space Models.} In recent times, State Space Models (SSMs) have garnered significant attention due to their linear scalability with respect to context length and competitive results relative to Transformers. Among the earliest architectures introduced was the Structured State Space sequence model (S4)~\citep{Gu:21}. This was succeeded by developments such as the Diagonal State Space (DSS) model~\citep{Gupta:22}, Gated State Space (GSS) models~\citep{Mehta:22}, S5 model~\citep{Smith:22}, Bidirectional Gated SSM (BiGS)~\citep{Wang:22}, the H3 model~\citep{Fu:23}, and Mamba~\citep{Gu:24arxiv}.

\paragraph{Recurrent Neural Networks.} In response to the challenge of achieving linear complexity with respect to context length, Recurrent Neural Networks (RNNs) have been proposed as alternatives to Transformer models and traditional attention mechanisms. Recent advancements include the use of Deep Linear Recurrent Units (LRUs) in RNNs, which have demonstrated strong performance in language modeling tasks~\citep{Orvieto:23, De:24arxiv}. Hierarchically Gated Linear RNN (HGRN)~\citep{Qin:23} and its successor HGRN2~\citep{Qin:24arxiv} also contribute notable improvements within this domain. Among notable large language modeling frameworks, RWKV~\citep{Peng:23arxivshort,Peng:24arxivshort} adopts an RNN-based approach and achieves results competitive with Transformer architectures.

\paragraph{Gating.}
The mechanism of gating, foundational to LSTM architectures, has been revisited and reinterpreted in a variety of recent models. Variants leveraging gating include HGRN~\citep{Qin:23}, HGRN2~\citep{Qin:24arxiv}, Gated Linear Attention (GLA)~\citep{Yang:23arxiv}, Gated State Space (GSS) models~\citep{Mehta:22}, Bidirectional Gated SSM (BiGS)~\citep{Wang:22}, Moving Average Equipped Gated Attention (MEGA)~\citep{Ma:22}, RWKV~\citep{Peng:23arxivshort}, and Mamba~\citep{Gu:24arxiv}.

\paragraph{Covariance Update Rule.} In order to improve storage capacity, the mLSTM cell was augmented with a matrix memory that operates according to a covariance update rule. Various existing approaches also rely on similar update mechanisms, such as Fast Weight Programmers \citep{Schmidhuber:92ncfastweights,Schlag:21}, RWKV-5 and RWKV-6~\citep{Peng:24arxivshort}, Retention~\citep{Sun:23arxiv}, Linear Transformer~\citep{Katharopoulos:20}, and HGRN2~\citep{Qin:24arxiv}.

\paragraph{Most Related.} In terms of conceptual similarity, the architectures most closely aligned with xLSTM include Retention~\citep{Sun:23arxiv}, RWKV~\citep{Peng:23arxivshort,Peng:24arxivshort}, and HGRN2~\citep{Qin:24arxiv}. These systems incorporate mechanisms such as matrix-based memory and gating functions. Nonetheless, unlike the recently introduced sLSTM, these methods do not support memory mixing capabilities. The inclusion of memory mixing is crucial for addressing state tracking tasks, thereby enhancing the expressive power of LSTMs relative to State Space Models (SSMs) and Transformers~\citep{Merrill:24,Deletang:23}. Such state tracking capacities are essential for tasks like code execution and maintaining continuity of entities throughout extended narratives.

\paragraph{Residually Stacking Architectures.} Similar to the architecture of nearly all state-of-the-art large-scale deep learning systems, xLSTM models are designed by arranging their building blocks in a residual, stacked manner~\citep{Srivastava:15,He:16}.

Such a design paradigm has paved the way for deep convolutional models~\citep{He:16} as well as the emergence of Transformers~\citep{Vaswani:17}. Transformers, in particular, have powered the rise of Large Language Models (LLMs) such as GPT-3~\citep{Brown:20short}, ChatGPT~\citep{Schulman:22short}, GPT-4~\citep{Achiam:23gpt4short}, Megatron-LM~\citep{Shoeybi:19}, Gopher~\citep{Rae:21short}, ERNIE 3.0 Titan~\citep{Wang:21arxivshort}, GLaM~\citep{Du:21glamshort}, Chinese M6~\citep{Lin:21}, multilingual AlexaTM 20B~\citep{Soltan:22}, OPT~\citep{Zhang:22opt}, Chinchilla~\citep{Hoffmann:22short}, BLOOM~\citep{Scao:22arxivshort}, GLM-130B~\citep{Zeng:22arxivshort}, LaMDA~\citep{Thoppilan:22short}, PaLM~\citep{Chowdhery:22short}, Llama~\citep{Touvron:23arxiv}, and Gemini~\citep{GeminiTeam:23arxiv,Reid:24arxiv}.

\section{Experiments}
\label{sec:Exp}

In this section, we present an empirical analysis of xLSTM, benchmarking it against state-of-the-art approaches, with a primary emphasis on language modeling. Section~\ref{sec:ExpSyntethic} explores the performance of xLSTM on a set of controlled, synthetic benchmarks to characterize its unique strengths. 
Section~\ref{sec:ExpComparison} provides a comparative study of validation set perplexities achieved by leading language modeling architectures, each trained on a 15B token subset of SlimPajama~\citep{Soboleva:23}. Within this experimental framework, we conduct ablation experiments to further dissect the workings of xLSTM, and investigate scaling trends following the methodologies described in \citet{Kaplan:20} and \citet{Brown:20short}. 
A more extensive language modeling analysis is performed in Section~\ref{sec:ExpLanguage}, 
where we evaluate xLSTM alongside top contenders from Section~\ref{sec:ExpComparison}, all trained on a significantly larger corpus of 300B SlimPajama tokens~\citep{Soboleva:23}. This section includes: (1) an assessment of the models' generalization to sequences of greater length, (2) comparison of validation perplexities and results on a suite of downstream tasks~\citep{Sutawika:24}, (3) systematic evaluation across the 571 domains of the PALOMA language benchmark suite~\citep{Magnusson:23arxivshort}, and (4) a scaling analysis—now utilizing a training set that is 20 times larger to elucidate efficiency and generalization properties across varying model and data scales.

\label{sec:xlstm_blocks_notation}
Throughout all experimental settings, we denote xLSTM[$a$:$b$] as representing a configuration with $a$ mLSTM-based blocks for every $b$ sLSTM-based blocks. As an illustration, xLSTM[7:1] corresponds to an arrangement of eight total blocks, among which seven employ mLSTM and one utilizes sLSTM. In scenarios where the total number of blocks is set to 48, this configuration yields 42 mLSTM-based blocks alongside 6 sLSTM-based blocks. Additionally, our experiments consistently employ both pre up-projection blocks for mLSTM and post up-projection blocks for sLSTM in all tested models.

\subsection{Synthetic Tasks and Long Range Arena}
\label{sec:ExpSyntethic}
We begin by evaluating xLSTM’s innovative exponential gating combined with memory mixing using formal languages as proposed in~\citep{Deletang:23}. Next, the utility of the novel matrix memory in xLSTM is measured on the Multi-Query Associative Recall task introduced by~\citep{Arora:23arxiv}. Lastly, we analyze xLSTM’s capabilities on long sequence processing through experiments conducted on the Long Range Arena benchmark~\citep{Tay:21}.

\paragraph{Evaluation of xLSTM's Exponential Gating Mechanism with Memory Mixing.} We assess the performance of xLSTM equipped with its novel exponential gating and memory mixing, mechanisms hypothesized to improve state tracking capabilities~\citep{Merrill:24, Merrill:23}. Our study builds on, and broadens, the suite of formal language benchmarks introduced by \citet{Deletang:23}, modifying them to incorporate training on sequences of varying lengths to allow for length extrapolation analyses. Comprehensive information regarding the task suite and an extended presentation of results is provided in Appendix~\ref{sec:appExpSynthetic-formalLang}. 
To benchmark xLSTM, we compare its performance with a range of alternative approaches, including Transformers, State Space Models, and other Recurrent Neural Networks. For evaluation, we focus on task-relevant tokens, measuring accuracy on a scale where 0 corresponds to random performance and 1 indicates perfect results. Specifically, we assess the following two-block variants across these tasks: xLSTM[0:1] (representing only sLSTM), xLSTM[1:0] (representing only mLSTM), xLSTM[1:1], Llama, Mamba, RWKV, Retention, Hyena, LSTM, as well as LSTM embedded within Transformer blocks (denoted as LSTM (Block)). Experimental outcomes are depicted in Figure~\ref{fig:formal-main}. 
Notably, models like Transformers and State Space Models that lack memory mixing and explicit state tracking, are incapable of solving even straightforward regular language tasks (such as parity). This aligns with previously reported evidence that both Transformers and State Space architectures are theoretically less expressive compared to RNNs~\citep{Merrill:24, Merrill:23, Deletang:23}.

\paragraph{Test of xLSTM's Memory Capacities on Associative Recall Tasks.} This experiment evaluates the novel matrix memory component of xLSTM by assessing its memory limits using the Multi-Query Associative Recall task~\citep{Arora:23arxiv}. In this setting, each input sequence contains key-value pairs sampled randomly from a large vocabulary, which must then be accurately stored and retrieved when queried. To further raise the task's complexity, we enlarge the number of key-value associations to as many as 256 and expand sequence context lengths to up to 2048 tokens. This permits a broader assessment of the memory capabilities across several model architectures. Our comparison covers 2-block variants of Llama, Mamba, RWKV-5, RWKV-6, xLSTM[1:1], and xLSTM[1:0], with performance measured by the proportion of correct recalls. Given that Transformer-based models such as Llama are known to possess exponential memory relative to coding dimension~\citep{Ramsauer:21}, they serve as an upper benchmark for this task. Figure~\ref{fig:mqar-main} displays the outcomes. 
Among models not based on the Transformer architecture, xLSTM[1:1] achieves the highest recall accuracy, including in smaller configurations. Notably, introducing the sLSTM block does not reduce overall memory capability; in fact, it enhances it, particularly observable in scenarios involving 256 key-value pairs. Further experiments, detailed in Appendix~\ref{sec:appExpSynthetic-mqar}, suggest through extrapolation analysis that xLSTM's improved memory enables successful generalization to contexts exceeding the lengths encountered in training.

\paragraph{Test of xLSTM's Long Context Capabilities on Long Range Arena.} To evaluate how xLSTM handles extended sequences and substantial context lengths, we benchmark it against alternative approaches using the Long Range Arena~\citep{Tay:21}. xLSTM consistently achieves high results across each of the tasks, indicating that the xLSTM model is highly adept at managing the challenges associated with long-context scenarios.

For more details, see Appendix~\ref{sec:appExpSynthetic-lra}.

\subsection{Method Comparison and Ablation Study}
\label{sec:ExpComparison}

In order to explore the central question of this study—namely, the capabilities of our novel LSTM variants when applied to large-scale language modeling—we conduct experiments training xLSTMs, Transformers, State Space Models, and additional approaches using 15B tokens sourced from SlimPajama, all within an auto-regressive training regime. Performance is evaluated by assessing these models on the validation dataset, and ablation experiments are conducted specifically for the xLSTM architectures.

\paragraph{Comparing xLSTM to Other Methods.} To provide a comprehensive comparison, we train various models on a dataset of 15B tokens from SlimPajama~\citep{Soboleva:23}, assessing their performance in terms of perplexity on a designated validation set. The set of evaluated architectures includes: xLSTM (our proposed approach), GPT-3 (Transformer)~\citep{Brown:20short}, Llama (Transformer)~\citep{Touvron:23arxiv}, H3 (SSM)~\citep{Fu:23}, Mamba (SSM)~\citep{Gu:24arxiv}, RWKV-4 (RNN)~\citep{Peng:23arxivshort}, RWKV-5 (RNN)~\citep{Peng:24arxivshort}, RWKV-6 (RNN)~\citep{Peng:24arxivshort}, GLA (linear Transformer)~\citep{Yang:23arxiv}, HGRN (RNN)~\citep{Qin:23}, HGRN2 (RNN)~\citep{Qin:24arxiv}, RetNet (linear Transformer)~\citep{Sun:23arxiv}, Hyena (linear Transformer)~\citep{Poli:23}, as well as variants xLSTM[1:0] and xLSTM[7:1] (see Section~\ref{sec:xlstm_blocks_notation}). Mixed-precision training was employed across the board; however, for RWKV-5, RWKV-6, GLA, and HGRN2, the mixed-precision approach did not involve PyTorch's automatic mixed-precision module (see Appendix Section~\ref{sec:appGenTrainProc} for further details).

We organize the models into three principal categories: (a) Transformers, (b) State Space Models (SSMs), and (c) Recurrent Neural Networks (RNNs), which include linear Transformers. Linear Transformers, in this context, refer to methods offering linear-time replacements for the traditional Transformer attention mechanism. For fairness, all models are sized comparably to a 350M parameter GPT-3 architecture, featuring a 1024-dimensional embedding and 24 residual layers. Notably, only GPT-3 implements weight sharing between token embeddings and output embeddings, resulting in a slightly reduced parameter count. Results reported in Table~\ref{tab:spaj15b_model_results} demonstrate that xLSTM achieves the lowest validation perplexity among all tested models. Further experimental specifics can be found in Appendix~\ref{sec:appExpComparison}. The scaling trends for this series of experiments, as depicted in Figure~\ref{fig:temp_sclaw15B}, suggest xLSTM's strong relative performance persists as model capacity increases.

\paragraph{Ablation Studies.}
As presented in Table~\ref{tab:spaj15b_model_results} and Figure~\ref{fig:temp_sclaw15B}, xLSTM demonstrates outstanding language modeling performance after being trained on 15B tokens from SlimPajama. To systematically investigate the impact of each architectural modification from LSTM to xLSTM, we progressively transform a standard LSTM architecture into xLSTM by making incremental adjustments. The first modification introduces LSTM layers into pre-LayerNorm residual frameworks. Next, the architecture is modified by incorporating a post up-projection block. The last step involves adding exponential gating along with matrix memory components.

The results are shown in Table~\ref{tab:ablstudies} (top). 

The ablation studies indicate that the notable enhancement in performance stems from the combined contributions of both the exponential gating and the matrix memory modules. Furthermore, given the critical role of gating within RNNs and State Space Models, we systematically investigate alternative gating strategies. As reported in Table~\ref{tab:ablstudies} (bottom), the results reveal that configuring each gate to be independently learnable and responsive to input yields consistent gains. Further detailed experiments focusing on the individual backbone components can be found in Appendix~\ref{sec:appAblDetails}.

\subsection{xLSTM as Large Language Model}
\label{sec:ExpLanguage}

This work concludes with comprehensive language modeling experiments on a large scale, evaluating xLSTM's capabilities as a large language model. To this end, we expand the dataset and train using 300B tokens drawn from SlimPajama. Notably, this token count matches that employed in prior work such as Mamba~\citep{Gu:24arxiv} and Griffin~\citep{De:24arxiv}.

We evaluate xLSTM in comparison to RWKV-4, Llama, and Mamba, each serving as a representative approach from the method categories outlined in Section~\ref{sec:ExpComparison}. RWKV-4 is chosen as the RNN exemplar because satisfactory training precision configurations for RWKV-5, RWKV-6, and HGRN2 became available only after initiating the 300B token training runs (refer to Appendix~\ref{sec:appGenTrainProc} for details). Multiple model scales are considered (125M, 350M, 760M, 1.3B parameters), with comprehensive testing for length extrapolation and validation performance. We further benchmark these models on a series of downstream tasks, assess their efficacy in language modeling across 571 text domains using the PALOMA benchmark, and examine their compliance with established scaling laws.

\paragraph{Sequence Length Extrapolation.}
\label{sec:spaj300B_ctx_extrapolation}
To begin, we evaluate how well 1.3B-parameter large-scale models—including xLSTM, RWKV-4, Llama, and Mamba—handle extrapolation to longer sequence lengths. Each model is trained with a fixed context length of 2048 tokens, then assessed on sequences extending up to 16384 tokens. Results are summarized in Figure~\ref{fig:spaj300B_seq_extrapolation}. Notably, xLSTM consistently achieves low perplexities as context length increases, outperforming other approaches on this metric.

\paragraph{Validation Perplexity and Downstream Tasks.}
\label{sec:spaj_downstream_eval}
Next, across every model scale, we assess xLSTM, RWKV-4, Llama, and Mamba on the SlimPajama validation suite for next-token prediction accuracy, as well as on downstream benchmarks focused on evaluating common sense reasoning capabilities.

The third column in Table~\ref{tab:lmeval} presents the validation set perplexities achieved by various approaches. For all examined model sizes, both xLSTM[1:0] and xLSTM[7:1] demonstrate superior performance regarding validation perplexity. The remaining columns in Table~\ref{tab:lmeval} display the results obtained on a range of downstream tasks. xLSTM consistently outperforms other methods across nearly all tasks and model sizes; exceptions occur only for the ARC task, where Mamba occasionally yields the best results. Further information is provided in Appendix~\ref{sec:appExpLanguage}.

\paragraph{Performance on PALOMA Language Tasks.} Third, we evaluate the next-token prediction abilities of xLSTM, RWKV-4, Llama, and Mamba models across all sizes using the PALOMA language benchmarks~\citep{Magnusson:23arxivshort}. We report perplexity scores for predicting subsequent tokens across 571 distinct text domains, spanning sources such as nytimes.com and Reddit's r/depression forum. Table~\ref{tab:ppleval} presents the resulting perplexity values, organized by language modeling (the first seven columns) and more specialized, domain-specific benchmarks (the final five columns). On these language tasks, xLSTM[1:0] demonstrates lower perplexity—and thus superior performance—compared to xLSTM[7:1].

xLSTM[1:0] achieves lower perplexity compared to Mamba in 568 of 571 text domains (99.5\%), outperforms Llama in 486 of 571 cases (85.1\%), and surpasses RWKV-4 in 570 of 571 domains (99.8\%). Additional information can be found in Appendix~\ref{sec:appExpLanguage}.

\paragraph{Scaling Laws.} Fourth, we investigate the power-law scaling patterns that enable predictions of model performance at larger scales~\citep{Kaplan:20,Brown:20short}. As shown in Figure~\ref{fig:temp_sclaw300B}, all architectures exhibit analogous scaling dynamics, although they differ in their baselines. RWKV-4 demonstrates the lowest performance, followed sequentially by Llama and Mamba. Notably, xLSTM surpasses Mamba, with the relative performance gap between xLSTM and Mamba mirroring that observed between Mamba and Llama. These scaling results suggest that, as models grow in size, xLSTM is expected to maintain a competitive edge relative to both Transformers and State-Space models.

\textbf{Generation Times and Maximal Throughput.} In Figures~\ref{fig:inference_time_generation} and~\ref{fig:inference_time_decoding_throughput} (left), we evaluate both the duration required to generate text and the maximum attainable throughput for our xLSTM variants at the 1.3B parameter level. We benchmark these results against Mamba, Llama, and RWKV implementations of similar model sizes available on HuggingFace, using a static key-value cache for Llama. During our evaluation, full cache compilation for Transformer models and Mamba compilation via \texttt{torch.compile} were not functional. All generation tests were conducted with a batch size of 1 and a pre-fill value of 16 for all models, a configuration that especially benefits the Transformer architecture.

As depicted in Figure~\ref{fig:inference_time_generation}, both xLSTM and the recurrent baselines (Mamba and RWKV-4) exhibit linear time complexity in text generation, in contrast to the quadratic scaling observed in the Llama model. Decoding throughput experiments further explore varying batch sizes and pre-fill values for Llama. Figure~\ref{fig:inference_time_decoding_throughput} (right) demonstrates that, owing to its constant memory footprint, xLSTM is able to handle much larger batch sizes than Llama and therefore attains superior throughput.

\section{Limitations}

(i) Unlike mLSTM, the sLSTM's memory mixing architecture precludes the possibility of executing operations in parallel, which in turn hampers fast parallel execution. Nonetheless, we engineered a custom CUDA kernel for sLSTM, resulting in a runtime that is under twice as slow as our parallel mLSTM version. (ii) The existing CUDA kernels used for mLSTM have not undergone optimization, leading to runtimes that are roughly four times slower compared to FlashAttention and the scan operation featured in Mamba. There is potential for significant acceleration by developing more efficient CUDA kernels analogous to those used for FlashAttention. (iii) mLSTM employs a matrix memory, introducing greater computational demands because the processing involves $d \times d$ matrices. However, memory read and write procedures are parameter-free and amenable to parallel execution with standard matrix routines; thus, the additional wall-clock time remains relatively insignificant. (iv) Special attention is required when selecting initial values for the forget gates. (v) Since matrix memory is decoupled from sequence length, handling substantially longer sequences could saturate the available memory, especially at increased context lengths. Despite this, no adverse effects have been observed for sequences up to 16k tokens, as discussed in Section~\ref{sec:spaj300B_ctx_extrapolation}. (vi) Owing to the high computational expense of scaling experiments to larger language models, we did not thoroughly tune the architecture or hyperparameters, particularly in larger xLSTM settings. Unlocking the full capabilities of xLSTM will necessitate more comprehensive optimization efforts.

\section{Conclusion}
Our investigation has partially addressed the central question: To what extent can language modeling benefit from scaling LSTM architectures to billions of parameters? Our current findings suggest that such scaling achieves performance at least on par with leading approaches such as Transformers and State Space Models. By introducing exponential gating combined with memory mixing and an updated memory architecture, we have advanced the standard LSTM to create the xLSTM variant. In empirical evaluations on language modeling tasks, xLSTM demonstrates strong performance relative to state-of-the-art techniques, including Transformers and State Space Models. Scaling law analyses suggest that expanding xLSTM models further could make them formidable alternatives to today's Large Language Models based on Transformer designs. Moreover, xLSTM holds promise for transformative advances in areas such as Reinforcement Learning, Time Series Prediction, and the modeling of physical systems.

\end{document}